\section{Related Work}
Various recently-proposed large-scale QA datasets can be categorized in four categories.
% have been constructed in recent years. We categorize them in four categories in our discussion.

\paragraph{Single-document datasets.} \squad{} \cite{rajpurkar2016squad, rajpurkar2018know} questions that are relatively simple because they usually require no more than one sentence in the paragraph to answer.
% collected with only one paragraph as the context.
% It is further constrained that the answer must appear as a span of text in the supporting paragraph, which enables evaluating QA systems' performance based on automatic metrics.
 % on their ability to match this span.
% This provides QA systems with stronger supervision, limits the search space, and allows for easier, more accurate automatic evaluation, which we also adopt in \datasetname{}.
% However,
% \squad{} questions are relatively simplistic because they usually require no more than one sentence in the paragraph to answer.
% Questions collected in this fashion are also less friendly to the open-context setting we are interested in, because they often don't contain sufficient information about the context.
% Finally, knowing what paragraph the answer is in is also unrealistic for real-world QA applications.

\paragraph{Multi-document datasets.} \triviaqa{} \cite{Joshi2017TriviaQA} and \searchqa{} \cite{dunn2017searchqa} contain question answer pairs that are accompanied with more than one document as the context.
This further challenges QA systems' ability to accommodate longer contexts.
However, since the supporting documents are collected after the question answer pairs with information retrieval, the questions are not guaranteed to involve interesting reasoning between multiple documents.
% Moreover, answers are induced in a distantly supervised manner, which does not guarantee the existence of an instance of the answer span that actually answers the question.

\paragraph{KB-based multi-hop datasets.}
Recent datasets like
\wikihop{} \cite{welbl2017constructing} and \complexwebq{} \cite{talmor2018web} explore different approaches of using pre-existing knowledge bases (KB) with pre-defined logic rules to generate valid QA pairs, to test QA models' capability of performing multi-hop reasoning.
% In each instance, questions are generated by taking a chain of entities or rules of reasoning in the KB or logical forms, and the answer is generated directly from the knowledge base.
% Although the QA system does not have access to the KB,
The diversity of questions and answers is largely limited by the fixed KB schemas or logical forms.
% Furthermore, even if the questions are naturalized by crowd workers, they would usually appear odd in a real-world setting if more than 2 hops are considered.
% Furthermore, they also fall prey to the incompleteness of the KBs -- if a fact is simply not populated into the knowledge base but is supported by textual evidence, they are at the risk of generating questions that seemingly require multi-hop reasoning, but are perfectly answerable with just one document.
Furthermore, some of the questions might be answerable by one text sentence due to the incompleteness of KBs.

\paragraph{Free-form answer-generation datasets.}
MS MARCO \cite{nguyen2016ms} contains 100k user queries from Bing Search with human generated answers. Systems generate free-form answers and are evaluated by automatic metrics such as ROUGE-L and BLEU-1. However, the reliability of these metrics is questionable because they have been shown to correlate poorly with human judgement \cite{novikova2017we}.

% 100k anonymized user queries from Bing Search over 1M passages over more than 200k documents, and contains human generated answers.
% Systems are allowed to freely generate answers given the context, and systems are evaluated by automatic metrics that focus mostly on $n$-gram overlap between the prediction and the gold reference, such as ROUGE-L and BLEU-1.
% This poses the problem of whether systems are being evaluated fairly and reliably, because these automatic metrics for automatic generated text have been shown to correlate poorly with human judgement \cite{}.
